%% ============================================================
%% Chapter 2 — Background and Related Work
%% ============================================================
\chapter{Background and Related Work}
\label{chap:background}

\section{Reading Comprehension Datasets}
\label{sec:bg_datasets}

\subsection{SQuAD 1.1}

The Stanford Question Answering Dataset (SQuAD) \cite{rajpurkar2016squad}
established a de-facto standard for machine reading comprehension when it was
released in 2016. The dataset consists of approximately 100,000 question–answer
pairs derived from 536 Wikipedia articles. Each question is paired with a
passage, and the answer is guaranteed to be a contiguous span within that
passage, making it amenable to extractive QA approaches. Human performance on
SQuAD 1.1 is approximately 82.3 F1 and 72.0 exact match, scores that modern
Transformer models have since surpassed.

A key limitation of SQuAD 1.1 is that \emph{every} question has an answer in
the provided passage. As a result, models trained on SQuAD 1.1 learn a
systematic bias toward always producing a span, regardless of whether the
passage actually contains an answer. This makes them unsuitable for open-domain
RAG settings where many queries are unanswerable given the retrieved context.

\subsection{SQuAD 2.0}

To address this limitation, \citet{rajpurkar2018squad2} released SQuAD~2.0
in 2018. The dataset augments SQuAD 1.1 with over 50,000 adversarially crafted
unanswerable questions, written by crowdworkers who were shown a paragraph and
asked to pose a question that looked plausible but could not be answered from
the paragraph. SQuAD~2.0 thus requires models to both locate answers in context
\emph{and} identify when no answer is present, returning a null span in the
latter case. Human performance stands at approximately 89.5 F1 and 86.9 exact
match, leaving a substantial gap to close.

The answerability challenge introduced by SQuAD~2.0 is directly aligned with
the abstention problem studied in this work. We use SQuAD~2.0 as the
experimental benchmark throughout.

\section{Sparse Retrieval with BM25}
\label{sec:bg_bm25}

Okapi BM25 (Best Match 25) \cite{robertson2009bm25} is the dominant sparse
retrieval function in information retrieval and has remained highly competitive
despite the rise of dense neural retrievers. Given a query $q$ and a document
(passage) $d$, BM25 scores the relevance of $d$ to $q$ as

\begin{equation}
  \text{BM25}(q, d) = \sum_{t \in q} \text{IDF}(t) \cdot
  \frac{f(t, d)\,(k_1 + 1)}
       {f(t, d) + k_1\!\left(1 - b + b\,\dfrac{|d|}{\text{avgdl}}\right)},
  \label{eq:bm25}
\end{equation}

\noindent where:

\begin{itemize}
  \item $f(t, d)$ is the term frequency of token $t$ in document $d$;
  \item $|d|$ is the length of document $d$ in tokens;
  \item $\text{avgdl}$ is the average document length in the collection;
  \item $k_1 \in [1.2, 2.0]$ controls term-frequency saturation;
  \item $b \in [0, 1]$ controls length normalisation (typically $b = 0.75$);
  \item $\text{IDF}(t)$ is the Inverse Document Frequency of token $t$,
        commonly computed as
        \[
          \text{IDF}(t) = \log\!\frac{N - n(t) + 0.5}{n(t) + 0.5},
        \]
        where $N$ is the total number of documents and $n(t)$ is the number of
        documents containing $t$.
\end{itemize}

Equation~\eqref{eq:bm25} has two notable properties: (i) term frequency
contributes with \emph{diminishing returns}, saturating as $f(t,d) \to \infty$;
and (ii) longer documents are \emph{penalised} through the length normalisation
factor $(1 - b + b\,|d|/\text{avgdl})$. These properties make BM25 robust
across heterogeneous corpora.

A key intuition motivating our BM25 Threshold strategy (Section~\ref{sec:meth_threshold})
is that the magnitude of the top-passage BM25 score encodes a form of
\emph{retrieval confidence}: a high score signals that the query terms overlap
substantially with the retrieved passage, making it more likely that the
passage contains the answer.

\section{Transformer-based Extractive QA}
\label{sec:bg_transformer_qa}

\subsection{BERT}

BERT (Bidirectional Encoder Representations from Transformers)
\cite{devlin2019bert} introduced a pre-trained Transformer encoder trained on
Masked Language Modelling (MLM) and Next Sentence Prediction (NSP) objectives.
When fine-tuned on SQuAD, BERT encodes the concatenated query and passage and
predicts two scalar logits per token—one for the start position and one for the
end position of the answer span. The token with the highest start logit and the
token with the highest end logit (after the start) jointly define the extracted
span.

\subsection{RoBERTa}

RoBERTa (Robustly Optimised BERT Pretraining Approach) \cite{liu2019roberta}
improved upon BERT through a more careful pre-training recipe: training on
larger data (160 GB vs 16 GB), longer sequences, larger batches, more steps,
dynamic masking, and removal of the NSP objective. These changes consistently
improved downstream task performance, including on SQuAD~2.0.

In this work we use the \texttt{deepset/roberta-base-squad2} checkpoint, a
RoBERTa-base model fine-tuned by the \texttt{deepset} team specifically on
SQuAD~2.0. This checkpoint is designed to handle the \texttt{handle\_impossible\_answer}
flag: it returns a \texttt{qa\_score} (probability that the extracted span is
the answer) and a \texttt{qa\_no\_answer\_probability} (probability that no
answer exists in the passage), both of which we use as features.

\section{Retrieval-Augmented Generation}
\label{sec:bg_rag}

\citet{lewis2020rag} introduced the RAG framework, which couples a dense
passage retriever (DPR) with a sequence-to-sequence generator. At inference
time, the retriever fetches the top-$k$ passages from a Wikipedia index; the
generator then conditions on the query and all retrieved passages to produce the
final answer. RAG demonstrated state-of-the-art results on open-domain QA
benchmarks and popularised the retriever–reader paradigm.

Subsequent work \cite{zhang2023rag} identified hallucination as a central
failure mode in RAG: the generator can produce confident-sounding text that is
not supported by any retrieved passage. This motivates the abstention paradigm
explored in our work, where the system learns to refrain from answering
precisely when the retrieval context is insufficient.

\section{Selective Prediction}
\label{sec:bg_selective}

Selective prediction \cite{geifman2017selective} is the framework in which a
model is allowed to \emph{abstain} on inputs where it is uncertain, rather than
always producing a prediction. Formally, a selective classifier is a pair
$(f, g)$ where $f$ is the prediction function and $g : \mathcal{X} \to \{0,1\}$
is the selection function; the system answers if and only if $g(x)=1$.
Performance is then measured over the \emph{selected} subset, together with the
\emph{coverage} $\Pr[g(x)=1]$.

A key result is that there exists a fundamental precision–coverage trade-off:
increasing coverage inevitably reduces precision at some point. The challenge is
to find a selection function that achieves a favourable operating point on this
frontier. In our work, the BM25 threshold and the logistic regression
classifier both serve as instantiations of $g$, with the threshold controlling
the trade-off.
